% ============================================================
%  CHAPTER 4 — METHODOLOGY (PROPOSAL)
%
%  Proposed reorganization (simulator-first, then RL, PLR last):
%
%   4.1  Overview                          (NEW — orients the reader)
%   4.2  The High-Throughput GPU Simulator (the enabling artifact)
%        4.2.1 Simulator state and step interface
%        4.2.2 Physics model
%        4.2.3 Batching and rollout collection
%   4.3  Problem Formulation               (centralized-joint-control MDP)
%        4.3.1 State and observation space
%        4.3.2 Action space
%        4.3.3 Reward function
%        4.3.4 Transition and discount
%   4.4  Policy Architecture               (centralized joint controller; CTDE deferred to TG-2)
%   4.5  Exploiting Symmetries
%   4.6  Two-Stage Curriculum and Self-Play
%   4.7  Prioritized-Level-Replay Sampler  (second contribution, own section)
%   4.8  Tooling and Infrastructure
%   4.9  Scope and Discarded Alternatives  (already written — keep at end)
%
%  Rationale: the simulator is the enabling artifact, so it comes first.
%  The MDP formulation then describes WHAT is being learned on top of it.
%  PLR is pulled out into its own section (not buried among RL details)
%  because it is the planned sample-selection mechanism for TG-2.
% ============================================================

This chapter defines the proposed methodology.  It begins with the
high-throughput simulator, the enabling artifact on which everything else
depends, then formulates the 3v3 task as a multi-agent reinforcement-learning
problem, describes the policy architecture and the mechanisms that exploit the
domain's symmetries, and presents the two-stage curriculum and
prioritized-level-replay sampler interface.  The chapter closes with the
tooling and a statement of scope.

% ==================================================================
\section{Overview}
% ==================================================================

The methodology realizes the two contributions of
Section~\ref{sec:contributions} within a common training loop: the simulator
generates batched rollouts, a \gls{ppo} learner consumes them, and the
initialization mechanism determines which situations to sample next.  The
remaining sections specify each component.

% ============================================================
%  TG-1 training-infrastructure diagram (revised).
%
%  Structural changes from the initial version:
%   – Learner and Policy merged into one SB3 PPO box.
%     The standard SB3 OnPolicyAlgorithm is a single synchronous
%     process; a two-box design implied an actor–learner split
%     (e.g. IMPALA) that does not exist in TG-1.
%   – JAX↔PyTorch device boundary made the visual focal point,
%     matching the throughput argument in Chapter 5.
%     All per-step obs/act data crosses this line every step.
%   – Opponent team explicitly shown in Environments.
%   – "On-device training loop" container removed: it spanned
%     both the JAX device and the PyTorch/CPU learner, making
%     "on-device" self-contradictory.
%
%  Required preamble (already present in tese.tex):
%      \usepackage{tikz}
%      \usetikzlibrary{arrows.meta,shapes,positioning}
% ============================================================

% ---- palette ----
\definecolor{tgDeepBlue}{RGB}{17,52,160}
\definecolor{tgTileBlue}{RGB}{96,141,224}
\definecolor{tgPaleBlue}{RGB}{223,231,249}
\definecolor{tgGold}{RGB}{245,199,38}
\definecolor{tgPaleGold}{RGB}{253,243,205}
\definecolor{tgCrimson}{RGB}{214,40,87}
\definecolor{tgArrow}{RGB}{40,40,55}

% ---- centred grid of small tiles (drawn from current origin) ----
% #1=lastCol (cols-1)  #2=lastRow (rows-1)  #3=colour  #4=step  #5=cell
\providecommand{\tgtiles}[5]{%
  \foreach \tx in {0,...,#1}{%
    \foreach \ty in {0,...,#2}{%
      \fill[#3] (\tx*#4 - #1*#4/2 - #5/2, \ty*#4 - #2*#4/2 - #5/2)
                rectangle +(#5,#5);}}%
}

\begin{figure}[ht]
    \centering
    \resizebox{\textwidth}{!}{%
    \begin{tikzpicture}[
        font=\scriptsize,
        card/.style={draw=tgDeepBlue, line width=0.8pt, rounded corners=2pt},
        flow/.style={-{Latex[length=2.2mm,width=1.8mm]}, line width=0.9pt, draw=tgArrow},
        lab/.style={font=\scriptsize, text=tgArrow, fill=white, inner sep=1.2pt},
        sub/.style={font=\tiny, align=center, text=tgDeepBlue!85!black},
    ]

    % ===== Background tints: CPU side (left) vs JAX side (right) =============
    % White strip at x=7.3–7.7 keeps the boundary line visually clear.
    \fill[tgPaleGold!30] (-0.2,-2.3) rectangle (7.3,5.7);
    \fill[tgPaleBlue!30] (7.7,-2.3) rectangle (15.2,5.7);

    % ===== JAX ↔ PyTorch device boundary — visual focal point ================
    \draw[line width=1.5pt, draw=tgCrimson] (7.5,-2.3) -- (7.5,5.7);
    \node[fill=white, text=tgCrimson, font=\tiny\bfseries,
          inner sep=2pt, align=center]
        at (7.5,5.5) {device\,$\leftrightarrow$\,host\\[-1pt](per-step)};

    % ===== Side labels =======================================================
    \node[font=\tiny\itshape, text=tgDeepBlue!70]
        at (3.0,-2.15) {CPU / PyTorch (host)};
    \node[font=\tiny\itshape, text=tgDeepBlue!70]
        at (12.2,-2.15) {JAX accelerator (device)};

    % ========================= SB3 PPO card (CPU side) =======================
    \node[card, fill=tgPaleGold, minimum width=5.2cm, minimum height=3.0cm]
        (sb3) at (3.0,0) {};
    \begin{scope}
      \clip[rounded corners=2pt] (sb3.south west) rectangle (sb3.north east);
      \fill[tgDeepBlue] ([yshift=-0.5cm]sb3.north west) rectangle (sb3.north east);
    \end{scope}
    \node[white, font=\scriptsize\bfseries]
        at ([yshift=-0.25cm]sb3.north) {SB3 PPO (PyTorch $\cdot$ CPU)};
    \begin{scope}[shift={(3.0,0.15)}]
        \tgtiles{3}{2}{tgGold}{0.26}{0.18}
    \end{scope}
    \node[sub] at (3.0,-0.82)
        {MLP $\pi_\theta$ (actor+critic trunk)\\
         single-process \texttt{OnPolicyAlgorithm}\\
         PPO gradient updates};

    % ========================= ENVIRONMENTS card (JAX side) ==================
    \node[card, fill=tgPaleBlue, minimum width=5.2cm, minimum height=3.0cm]
        (env) at (12.2,0) {};
    \begin{scope}
      \clip[rounded corners=2pt] (env.south west) rectangle (env.north east);
      \fill[tgDeepBlue] ([yshift=-0.5cm]env.north west) rectangle (env.north east);
    \end{scope}
    \node[white, font=\scriptsize\bfseries]
        at ([yshift=-0.25cm]env.north) {Environments};
    \begin{scope}[shift={(12.2,0.15)}]
        \tgtiles{5}{3}{tgTileBlue}{0.22}{0.15}
    \end{scope}
    \node[sub] at (12.2,-0.82)
        {\texttt{VSSVecEnv} $\cdot$ JAX $\cdot$
         $\mathrm{jit}(\mathrm{vmap}(\texttt{step}))$\\
         256--1024 parallel matches\\
         opponent: scripted / frozen $\pi^{(t-k)}_\theta$};

    % ============================ MLflow card ================================
    \node[card, fill=tgPaleBlue, minimum width=3.8cm, minimum height=1.55cm]
        (mlflow) at (3.0,4.4) {};
    \begin{scope}
      \clip[rounded corners=2pt] (mlflow.south west) rectangle (mlflow.north east);
      \fill[tgDeepBlue] ([yshift=-0.42cm]mlflow.north west) rectangle (mlflow.north east);
    \end{scope}
    \node[white, font=\scriptsize\bfseries]
        at ([yshift=-0.21cm]mlflow.north) {MLflow tracking};
    \begin{scope}[shift={(3.0,4.55)}]
        \tgtiles{4}{1}{tgTileBlue}{0.2}{0.13}
    \end{scope}
    \node[sub] at (3.0,4.12)
        {reward (mean/min) $\cdot$ ep.\ length\\throughput $\cdot$ checkpoints};

    % ========================= INITIALIZATION card ===========================
    \node[card, fill=tgPaleBlue, minimum width=4.2cm, minimum height=1.55cm]
        (init) at (12.2,4.4) {};
    \begin{scope}
      \clip[rounded corners=2pt] (init.south west) rectangle (init.north east);
      \fill[tgDeepBlue] ([yshift=-0.42cm]init.north west) rectangle (init.north east);
    \end{scope}
    \node[white, font=\scriptsize\bfseries]
        at ([yshift=-0.21cm]init.north) {Initialization};
    \begin{scope}[shift={(12.2,4.55)}]
        \tgtiles{4}{1}{tgTileBlue}{0.2}{0.13}
    \end{scope}
    \node[sub] at (12.2,4.12)
        {\texttt{InitMode}: kickoff / random\\($\rightarrow$ PLR sampler, TG-2)};

    % ============================ ARROWS =====================================

    % Per-step: allied action SB3 → Env (crosses device boundary)
    \draw[flow] ([yshift=0.5cm]sb3.east) -- ([yshift=0.5cm]env.west)
        node[lab, pos=0.25] {act $a\!\in\!\mathbb{R}^{6}$ (allied)};
    % Per-step: observation Env → SB3 (crosses device boundary)
    \draw[flow] ([yshift=-0.5cm]env.west) -- ([yshift=-0.5cm]sb3.east)
        node[lab, pos=0.75] {obs $o\!\in\!\mathbb{R}^{46}$};

    % Gymnasium ⇌ SB3 adapter sitting ON the boundary —
    % mediates every per-step obs/act transfer
    \node[draw=tgDeepBlue, line width=0.8pt, fill=white, rounded corners=2pt,
          align=center, font=\tiny, inner sep=2.5pt]
        at (7.5,0) {Gymnasium\,$\rightleftarrows$\,SB3\\adapter};

    % Rollout buffer returned after N steps (bottom path)
    \draw[flow] (env.south) -- (12.2,-1.85) -- (3.0,-1.85) -- (sb3.south);
    \node[lab] at (7.5,-1.85)
        {rollout $(o,a,r,\text{done})$ — $N$ steps};

    % Initialization → Environments
    \draw[flow] (init.south) -- (env.north)
        node[lab, midway] {initial states};

    % SB3 → MLflow
    \draw[flow] (sb3.north) -- (mlflow.south)
        node[lab, midway] {metrics};

    \end{tikzpicture}%
    }
    \caption{Training infrastructure for the TG-1 smoke tests
    (\texttt{scripts/smoke.py}).  The figure is organised around the
    \textbf{device\,$\leftrightarrow$\,host boundary} (red vertical line):
    the CPU/PyTorch side (left, warm tint) holds the PPO learner, and the
    JAX accelerator side (right, blue tint) holds the vectorised
    environments.  \emph{Every step}, the environment receives the allied
    action $a\!\in\!\mathbb{R}^{6}$ (wheel speeds for all three robots
    output by the centralized MLP $\pi_\theta$) and an opponent action
    (scripted agent or frozen checkpoint $\pi^{(t-k)}_\theta$), then
    returns observation $o\!\in\!\mathbb{R}^{46}$ (positions, velocities,
    and headings of all six robots plus the ball).  Both arrays must cross
    the Gymnasium\,$\rightleftarrows$\,SB3 adapter on every call
    (device$\to$host for obs; host$\to$device for act); this per-step
    transfer is the throughput bottleneck motivating the JAX-native
    training loop planned for TG-2.  The \textbf{SB3 PPO box is a single
    synchronous process} (\texttt{OnPolicyAlgorithm}): there is no
    actor--learner split; \texttt{collect\_rollouts} and \texttt{train}
    share the same MLP object (actor and critic share an MLP trunk ---
    not per-agent parameter sharing); the learner is intentionally placed
    on CPU to avoid GPU contention with the JAX device.  After $N$ steps
    the accumulated rollout buffer feeds the PPO update.  Episode resets
    use the pluggable \texttt{InitMode} (the future PLR sampler insertion
    point), and MLflow logs per-episode mean/min reward, episode length,
    throughput, and checkpoints.  The centralized joint-control
    formulation (one MLP, one team) is the TG-1 baseline;
    per-agent parameter sharing and \gls{ctde} are deferred to TG-2.}
    \label{fig:training_infra}
\end{figure}



% ==================================================================
\section{The High-Throughput GPU Simulator}
% ==================================================================

The simulator is implemented as a pure, batched JAX environment, designed for
speed, numerical stability, and full on-accelerator execution so that the
entire collect-and-update loop runs without per-step transfers to host memory.

% ------------------------------------------------------------------
\subsection{Simulator State and Step Interface}
% ------------------------------------------------------------------

A simulator state stores the poses and velocities of the six robots, the
position and velocity of the ball, the score, episode timers, and the
pseudo-random generator keys.  In the JAX backend this state is an immutable
\texttt{NamedTuple} PyTree, which makes it directly compatible with
\texttt{jax.vmap} batching.  The environment exposes two functions:
%
\begin{itemize}
    \item \texttt{reset(key)} initializes one or more game states according to
    the selected initialization mode;
    \item \texttt{step(state, action)} advances the physics by one timestep and
    returns the next state together with observations, rewards, termination
    flags, and diagnostic information.
\end{itemize}

% ------------------------------------------------------------------
\subsection{Physics Model}
\label{sec:physics_model}
% ------------------------------------------------------------------

The physics model is two-dimensional and captures the interactions that shape
policy learning: differential-drive motion, ball dynamics under friction,
robot–ball and robot–robot contact, and wall and goal handling on an IEEE VSSS
field.  Numeric dimensions (field size, robot footprint, ball diameter, goal
opening, wheelbase, and control rate) follow the competition specification and
are centralized in the implementation configuration module
(\texttt{vsss\_sim.config}).  The subsections below state \emph{why} each
modelling choice was made; run-specific training overrides are reported
separately in Chapter~5 (Table~\ref{tab:smoke_config}).

\subsubsection{Scope and fidelity}
\label{sec:physics_scope}

TG-1 prioritizes a simulator that is fast, batched, and stable enough for
large-scale \gls{ppo} rollouts over reproducing every effect of the team's
legacy C++/ODE stack.  The model is planar: no vertical dynamics,
no wheel slip model, and no explicit motor controller beyond rate-limited wheel
commands.  That scope matches the scale of Very Small Size Soccer, where the
dominant learning signal comes from ball interaction and team positioning
rather than from high-fidelity drivetrain modeling.

\subsubsection{Geometry, coordinates, and goals}

The regulation field layout is Figure~\ref{fig:field} (Chapter~2).  In the
simulator the origin is at the field centre, $+x$ points toward the opponent
goal, and $+y$ is parallel to the goal line (Section~\ref{sec:ma_coords}).
Each team fields three robots as $7.5 \times 7.5$\,cm \textbf{oriented bounding
boxes (OBBs)} with circle--OBB ball contact (not circular robot proxies).
Goals are $40$\,cm openings; scoring triggers an in-episode kickoff, not a
terminal state.

\subsubsection{Time stepping and actuation}

Control commands arrive at $60$\,Hz ($\Delta t = 1/60$\,s).  Each control step
is subdivided into \textbf{four physics sub-steps} (\texttt{SUB\_STEPS}), so
collisions and ball motion are integrated at $240$\,Hz effective resolution
without increasing the policy's decision frequency.  Robots are
\textbf{differential-drive}: normalized actions in $[-1,1]$ map to left and
right wheel speeds, then to body-frame $(v_x, v_y, \omega)$ via the standard
kinematic relation with wheelbase $53$\,mm and maximum linear speed
$\sim 1.30$\,m/s.

Wheel speeds do not jump instantaneously to the commanded value.  Each sub-step
limits the change in wheel speed to approximate motor slew (\texttt{ROBOT\_WHEEL\_ACCEL\_LIMIT} $\approx 1\,g$), which smooths trajectories and avoids
unphysically sharp turns that would otherwise appear when a neural policy
outputs discontinuous commands.  Heading is wrapped to $(-\pi, \pi]$ after
integration.

\subsubsection{Contact dynamics and ball motion}

Between commands, the ball experiences rolling friction (coefficient
\texttt{BALL\_FRICTION}) and elastic-style impulses on walls and robots
(restititution parameters in \texttt{config}).  Robot–wall contact clamps the
OBB inside the field and zeroes velocity components that point into the wall.
Robot–robot collisions use separating axes on the square footprints with
impulse exchange.  These rules are hand-tuned for stable, non-exploding
contacts at VSSS speeds rather than for exact material parameters; the goal is
repeatable interaction statistics for learning, not a calibrated multibody
model.

\subsubsection{Physics design choices}
\label{sec:physics_design}

Table~\ref{tab:physics_design} summarizes the main modelling decisions and
their rationale.  Together they define the physics contract that both backends
implement.

\begin{table}[ht]
    \centering
    \caption{Physics design choices in the VSSS simulator.
    Numeric constants live in \texttt{vsss\_sim.config}.}
    \label{tab:physics_design}
    \small
    \begin{tabular}{p{0.22\linewidth}p{0.38\linewidth}p{0.34\linewidth}}
        \toprule
        \textbf{Choice} & \textbf{Mechanism} & \textbf{Rationale} \\
        \midrule
        Planar 2D engine &
        State: ball $(x,y,v_x,v_y)$; robots $(x,y,\theta,v_x,v_y,\omega)$ &
        Matches top-down VSSS; enables batched array physics. \\
        \addlinespace
        OBB robots &
        Square footprint; circle–OBB for ball; SAT for robot–robot &
        Avoids disk proxy errors at $7.5$\,cm scale. \\
        \addlinespace
        Sub-stepped integration &
        $4\times$ substeps per $60$\,Hz control step &
        Reduces tunneling and stabilizes contacts at max wheel speed. \\
        \addlinespace
        Rate-limited wheels &
        Slew cap $\approx 1\,g$ per sub-step on each wheel &
        Models torque limits; smoother policy rollouts. \\
        \addlinespace
        Differential drive &
        Wheel speeds $\rightarrow$ $(v,\omega)$ $\rightarrow$ world velocity &
        Matches IEEE VSS hardware; action space is six wheel commands. \\
        \addlinespace
        Rolling ball friction &
        Deceleration $\propto \mu_r g \,\Delta t_{\mathrm{sub}}$ &
        Ball comes to rest without explicit spin state. \\
        \addlinespace
        Restitution on impact &
        Separate coeffs.\ for ball–wall, ball–robot, robot–wall &
        Tunable bounce; prevents energy blow-up in dense clusters. \\
        \addlinespace
        Goals without terminal &
        Score event + kickoff reset; episode ends on step cap only &
        Aligns with continuous play; keeps \texttt{terminated}$=$\texttt{False}. \\
        \bottomrule
    \end{tabular}
\end{table}

% ------------------------------------------------------------------
\subsection{Batching and Rollout Collection}
% ------------------------------------------------------------------

The environment is vectorized across many parallel games by applying
$\mathrm{jit}(\mathrm{vmap}(\texttt{step}))$ to a single batched state, so that
one compiled call advances an entire batch of independent matches.  Rollouts for
\gls{ppo} are collected by scanning the step function through time and mapping
it across environments, keeping trajectory buffers resident on the device.

%% FIGURE: on-device training loop — batched SimState → jit(vmap(step)) →
%% observations/rewards → PPO update → next initial states (from the sampler).
%% Place beside the Chapter 2 socket-loop figure so the architectural shift is
%% visible; both should be your own diagrams in a consistent visual style.

% ==================================================================
\section{Problem Formulation}
% ==================================================================

The 3v3 \gls{vsss} task is the domain within which both contributions are
evaluated.  In TG-1, the team is controlled by a \textbf{centralized joint
policy}: a single network observes the full global state of all six robots and
the ball ($\mathbb{R}^{46}$) and outputs a joint action that simultaneously
commands all three allied robots ($\mathbb{R}^{6}$).  Although the domain is
inherently multi-agent, this centralized formulation treats the three-robot
team as a single composite agent, which simplifies the implementation and
provides a clean baseline.  The observation interface is factored as a
swappable function, so the per-robot egocentric formulation required by
\gls{ctde} can be adopted without changes to the simulator.  This section
defines the components of the underlying decision process: state and
observation space, action space, reward, and transition dynamics.

\label{sec:ma_coords}
The learning frame uses the centred simulator coordinates of
Figure~\ref{fig:field}: origin at field centre, $+x$ toward the opponent goal,
$+y$ along the goal line.

% ------------------------------------------------------------------
\subsection{State and Observation Space}
% ------------------------------------------------------------------

The full simulator state is described in Section~4.2.  From it, the learning
system builds a per-step observation vector of dimension 46, composed of the
ball state and the state of all six robots, normalized so that equivalent
situations have similar numerical representations.  The ball contributes its
position and velocity; each robot contributes its position, orientation, and
linear and angular velocity.  Orientation is encoded as the pair
$(\sin\theta, \cos\theta)$ rather than the raw angle, to avoid the
discontinuity at $\pm\pi$ and to give the network a smooth representation of
heading.  Table~\ref{tab:obs} summarizes the observation layout and the
normalization applied to each group.

\begin{table}[ht]
    \centering
    \caption{Observation vector ($\mathbb{R}^{46}$) layout and normalization.
    Positions are normalized by the half-field dimensions, linear velocities by
    $1.5\times$ the maximum wheel speed, and angular velocity by the maximum
    achievable rate implied by the wheelbase.}
    \label{tab:obs}
    \begin{tabular}{lll}
        \toprule
        \textbf{Indices} & \textbf{Content} & \textbf{Per element} \\
        \midrule
        $[0\!:\!4]$  & ball               & $x, y, v_x, v_y$ (normalized) \\
        $[4\!:\!46]$ & 6 robots $\times$ 7 & $x, y, \sin\theta, \cos\theta, v_x, v_y, \omega$ \\
        \bottomrule
    \end{tabular}
\end{table}

The observation construction is factored as a swappable function.  This keeps
the interface forward-compatible with per-robot egocentric observations, which
the multi-agent variants planned for TG-2 require, without changing the
simulator.

% ------------------------------------------------------------------
\subsection{Action Space}
% ------------------------------------------------------------------

The action space is continuous.  In the baseline formulation, the policy
controls all three allied robots through a single 6-dimensional action — a
normalized left/right wheel-speed pair per robot, in $[-1, 1]$ — which is
denormalized to physical units inside the backend.
The policy controls wheel speeds directly rather than emitting a higher-level
target position handed to an external controller.  This merges the
path-planning and control problems into a single learned policy: it raises task
difficulty but removes the performance ceiling that a fixed external controller
would impose, an effect also observed in earlier VSSS control
work~\parencite{tg_thiago, brandao2022}.

A body-frame command $(v, \omega)$ converted to wheel speeds through the
differential-drive model (Chapter~2) is an alternative interface that can make
exploration smoother.  The methodology keeps both interfaces explicit so that
the final implementation can compare them empirically.

% ------------------------------------------------------------------
\subsection{Reward Function}
% ------------------------------------------------------------------

The reward is team-based and shared across the three allied robots, since they
pursue a common objective.  At its core it is sparse:
%
\begin{equation}
    r_t = g_t + \beta\,\Delta x^{b}_t,
    \label{eq:reward}
\end{equation}
%
where $g_t \in \{+1, -1, 0\}$ is the goal signal (scoring, conceding, or
neither), $\Delta x^{b}_t$ is the forward progress of the ball toward the
opponent goal during the timestep, and $\beta$ is a small shaping coefficient.
The shaping term scaffolds early learning, when goal events are too rare to
provide useful gradient under a purely random policy.  The coefficient is
chosen so that the maximum shaping contribution accumulated over an episode is
strictly smaller than a single goal event; with $\beta = 0.10$ and a field
length of $1.50\,\text{m}$, the total dense reward is bounded by
$0.10 \times 1.50 = 0.15 < 1$, so the agent is never incentivized to shuttle
the ball indefinitely instead of scoring.  Sharing a single team reward keeps
credit assignment simple at the cost of making the individual contribution of
each robot harder to isolate; additional shaping terms for possession,
positioning, and avoidance of self-blocking are left as tuning options for
TG-2.

% ------------------------------------------------------------------
\subsection{Transition and Discount}
% ------------------------------------------------------------------

The transition dynamics are encoded entirely in the simulator and are not made
explicitly available to the agent.  The discount factor $\gamma \in [0, 1)$
weights future reward against immediate reward; its value is treated as a
training hyperparameter and reported with the results.

% ==================================================================
\section{Policy Architecture}
% ==================================================================

The TG-1 policy is a \textbf{centralized joint controller}: a single
\gls{mlp} receives the global $\mathbb{R}^{46}$ observation vector and
produces the full $\mathbb{R}^{6}$ joint action — left and right wheel speeds
for all three allied robots — in a single forward pass.  Both the actor and
the critic share this global view.  There is no per-robot decomposition of
the observation, no separate per-robot action head, and no decentralized
execution at test time.

This centralized formulation maps directly onto a standard single-agent
\gls{ppo} implementation and eliminates credit-assignment complexity, so that
the throughput and sampler contributions can be evaluated in isolation.  Its
main limitation is that the six-dimensional joint action conditions all robot
commands on a single global representation; a per-robot formulation with
egocentric observations and \gls{ctde} removes this constraint while keeping
execution decentralized.  The observation construction is already factored as
a swappable function so that this transition does not require any change to
the simulator.

% ==================================================================
\section{Exploiting Symmetries}
% ==================================================================

Even within the centralized formulation, two classes of symmetry of the
\gls{vsss} domain can be exploited to improve sample efficiency.

\textbf{Input-ordering invariance.}  The 46-dimensional observation concatenates
ball features followed by six robots each contributing seven features.  Because
the three allied robots are homogeneous, any permutation of their slots in the
input produces an equivalent game state; ideally the policy should be invariant
to this relabeling.  In TG-1 the robot ordering in the observation is fixed and
consistent across steps, so the network must learn this invariance from data.
A permutation-equivariant architecture, or a per-robot formulation that makes
the invariance exact by construction, eliminates this dependency on learned
invariance.

\textbf{Reflection augmentation.}  The \gls{vsss} field is bilaterally symmetric
about its longitudinal axis.  A transition $(o_t, a_t, r_t, o_{t+1})$ collected
from play can be mirrored — transverse coordinates and velocities negated,
transverse components of the joint action negated — to produce a second valid
transition at no additional simulation cost.  This augmentation applies directly
to the centralized formulation and is expected to double the effective number of
distinct states the policy is trained on, reducing overfitting to specific field
positions.  The details of the sign-flip transform for positions, velocities, and
wheel-speed actions are implemented in the simulator's observation module.

% ==================================================================
\section{Curriculum Learning and Self-Play}
% ==================================================================

Training employs a curriculum that gradually increases the difficulty and complexity of the environment as learning progresses, rather than exposing the agent to the full challenge from the outset. An initial phase may involve only the allied robots and the ball, allowing the policy to develop basic ball handling skills such as reaching, conducting, and coordinated positioning without interference. As proficiency improves, one or more opponent robots are introduced, either as stationary obstacles, simple scripted agents, or through self-play.

In self-play regimes, adversaries are sampled from current or earlier versions of the policy, ensuring that the policy faced by the allies is both challenging and representative of varied behaviors.  Opponent sampling strategies — such as always using the latest policy or mixing against a population weighted by recency or performance~\cite{master_thay,starcraft_marl} — trade off training stability against generalization to diverse opponent styles.

Advancement through the curriculum can be based on various criteria, such as a fixed number of steps, reaching performance thresholds, or a combination of both. The concrete progression mechanism may be tuned to optimize training stability and efficiency, and will be finalized in subsequent development stages.

% ==================================================================
\section{Prioritized-Level-Replay Sampler}
% ==================================================================

Uniform episode initialization spends many simulator steps on situations that
are either trivially solved or too chaotic to yield useful gradients.  The
proposed sampler treats each initialized game situation as a replayable level
and samples levels according to their estimated learning potential.  Together,
the simulator and the sampler address orthogonal costs: the former increases
raw transition throughput, the latter improves the informativeness of each
transition.

The sampler maintains a population of candidate \gls{vsss} scenarios, each
defined by robot poses, ball state, velocities, and optional game context.
During training, each scenario receives a priority derived from learning
signals — primarily the magnitude of the value-prediction error on recent
trajectories, optionally combined with policy entropy or recent improvement.
Episodes are then initialized from a mixture of high-priority scenarios and
exploratory samples, which prevents the agent from overfitting to a narrow set
of situations.  The infrastructure for this already exists in the simulator:
the pluggable initialization interface (Chapter~5) was designed to accept an
external placement source, which is exactly the insertion point the sampler
requires.

For TG-1, the contribution is the sampler design and its integration with the
simulator and the \gls{ppo} rollout interface.  TG-2 will evaluate whether the
sampler improves sample efficiency relative to uniform initialization and to
hand-designed curricula.

% ==================================================================
\section{Tooling and Infrastructure}
% ==================================================================

The implementation uses JAX for numerical computation and accelerator
execution.  The reinforcement-learning stack is centered on \gls{ppo}, with
rollout collection and policy updates designed to minimize transfers between
host and device memory.  Experiment tracking records simulator throughput,
training curves, reward components, environment counts, hardware information,
and policy checkpoints.  The primary engineering requirement for TG-1 is that
the simulator be \emph{measurable}: throughput is reported as steps per second
versus the number of parallel environments, and compared against a vectorized
\gls{cpu} baseline, in Chapter~5.

% ==================================================================
\section{Scope and Discarded Alternatives}
% ==================================================================

The methodology reflects deliberate scoping decisions. Several reasonable
directions were excluded so the work could concentrate on the accelerated
simulator, the centralized training baseline, and the sample-selection path they
enable.

\subsection{Task-specific policies.}
Much prior \gls{vsss} work trains policies for isolated sub-tasks (interception,
penalty kicks), which are easier to learn and evaluate. This work instead targets
a single policy for both offense and defense across a full 3v3 match. Team
coordination only emerges when the policy owns the complete game, and the
resulting difficulty is exactly what the simulator and sample-selection
mechanism are meant to absorb.

\subsection{Alternative learning algorithms.}
\gls{ppo} was fixed for comparability, not superiority: the recent ITAndroids
results \parencite{medeiros2020, tg_thay, brandao2022} all use it, keeping this
work comparable with the team's baselines. Since the contributions are orthogonal
to the algorithm, as a faster simulator and better initialization strategy benefit
any on-policy method, fixing \gls{ppo} removes a degree of freedom without
limiting generality.

\subsection{Existing GPU RL frameworks.}
General accelerated simulators (Isaac Gym, Brax) target full 3D rigid-body
dynamics, far heavier than \gls{vsss} needs. The domain is simple enough — 2D,
differential drive, OBB collisions, static walls — that a purpose-built batched
kernel achieves higher throughput precisely because it computes less.

\subsection{Wrapping the existing C++/ODE simulator.}
The existing ODE simulators (Chapter~2) are bottlenecked architecturally —
single-instance, CPU-bound, socket-mediated — not by tuning. No wrapper makes a
single-threaded engine step thousands of environments on-device, so reuse would
preserve the exact limitation the new simulator removes.

\subsection{Domain randomization.}
Varying friction, latency, and noise for transfer robustness is sequenced rather
than discarded: it serves neither throughput nor sample efficiency, and becomes
relevant only once a competitive policy exists and the goal shifts to physical
deployment (Chapter~6).